% =========================================================
% Section: Sigma-Algebras
% Chapter: Algebras of Sets — Volume I
% =========================================================

\begin{tcolorbox}[colback=gray!6, colframe=gray!40, arc=2pt,
  left=6pt, right=6pt, top=4pt, bottom=4pt,
  title={\small\textbf{Section Toolkit --- Quick Reference}},
  fonttitle=\small\bfseries]
\small
\begin{tabular}{l l l l}
\toprule
\textbf{Label} & \textbf{Statement} & \textbf{Proof method} & \textbf{Detail} \\
\midrule
D\ref{def:generated-sigma-algebra}
  & $\sigma(\mathcal{C})$ defined as smallest $\sigma$-algebra
    containing $\mathcal{C}$
  & Definition & \hyperref[def:generated-sigma-algebra]{$\downarrow$} \\
P\ref{prop:generated-sigma-exists}
  & $\sigma(\mathcal{C})$ exists
  & $\mathcal{P}(X)$ witnesses non-emptiness; intersect
  & \hyperref[prf:generated-sigma-exists]{$\downarrow$} \\
P\ref{prop:generated-sigma-idempotent}
  & $\sigma(\sigma(\mathcal{C})) = \sigma(\mathcal{C})$
  & Minimality argument
  & \hyperref[prf:generated-sigma-idempotent]{$\downarrow$} \\
P\ref{prop:generated-sigma-monotone}
  & $\mathcal{C}_1\subseteq\mathcal{C}_2
    \Rightarrow\sigma(\mathcal{C}_1)\subseteq\sigma(\mathcal{C}_2)$
  & Minimality argument
  & \hyperref[prf:generated-sigma-monotone]{$\downarrow$} \\
P\ref{prop:countable-cocountable-sigma}
  & Countable/co-countable sets form a $\sigma$-algebra
  & Verify axioms directly
  & \hyperref[prf:countable-cocountable-sigma]{$\downarrow$} \\
P\ref{prop:finite-cocountable-algebra-not-sigma}
  & Finite/co-finite sets on $\mathbb{N}$ form an algebra
    but not a $\sigma$-algebra
  & Explicit construction of a witness
  & \hyperref[prf:finite-cocountable-algebra-not-sigma]{$\downarrow$} \\
P\ref{prop:sigma-from-singletons}
  & On a countable set, $\sigma(\{\{x\}: x\in X\}) = \mathcal{P}(X)$
  & Every subset is a countable union of singletons
  & \hyperref[prf:sigma-from-singletons]{$\downarrow$} \\
P\ref{prop:sigma-closed-symmdiff}
  & $\sigma$-algebra closed under countable symmetric difference
  & Express $\triangle$ via $\cup$, $\cap$, $(\cdot)^c$
  & \hyperref[prf:sigma-closed-symmdiff]{$\downarrow$} \\
\bottomrule
\end{tabular}
\end{tcolorbox}

% ---------------------------------------------------------
% Generated sigma-algebra definition
% ---------------------------------------------------------

\begin{tcolorbox}[colback=propbox, colframe=propborder, arc=2pt,
  left=6pt, right=6pt, top=4pt, bottom=4pt,
  title={\small\textbf{Definition (Generated $\sigma$-Algebra)}},
  fonttitle=\small\bfseries]
\label{def:generated-sigma-algebra}
Let $\mathcal{C} \subseteq \mathcal{P}(X)$.  The
\textbf{$\sigma$-algebra generated by $\mathcal{C}$}, denoted
$\sigma(\mathcal{C})$, is defined as
\[
  \sigma(\mathcal{C})
  := \bigcap \bigl\{
       \mathcal{F} \subseteq \mathcal{P}(X) :
       \mathcal{F} \text{ is a } \sigma\text{-algebra on } X
       \text{ and } \mathcal{C} \subseteq \mathcal{F}
     \bigr\}.
\]
\end{tcolorbox}

\begin{remark*}[Fully quantified form]
$\forall\mathcal{C}\subseteq\mathcal{P}(X)$:
$\sigma(\mathcal{C})$ is the intersection of all $\sigma$-algebras
on $X$ that contain $\mathcal{C}$.
\end{remark*}

\begin{remark*}[Flash]
\FlashQ{How is the generated $\sigma$-algebra $\sigma(\mathcal{C})$ defined?}
\FlashA{$\sigma(\mathcal{C}) := \bigcap\{\mathcal{F} :
  \mathcal{F}\ \sigma\text{-algebra on }X,\ \mathcal{C}\subseteq\mathcal{F}\}$.
  It is the smallest $\sigma$-algebra on $X$ containing $\mathcal{C}$.}
\FlashHint{The definition is formal; the content is in the existence proof,
  which requires showing the intersecting family is non-empty.}
\end{remark*}

% ---------------------------------------------------------
% Propositions
% ---------------------------------------------------------

\begin{proposition}[\texorpdfstring{%
  \hyperref[prf:generated-sigma-exists]{%
    Generated $\sigma$-Algebra Exists}}{%
  Generated Sigma-Algebra Exists}]
\label{prop:generated-sigma-exists}
For any $\mathcal{C} \subseteq \mathcal{P}(X)$, $\sigma(\mathcal{C})$
exists and is the unique smallest $\sigma$-algebra on $X$ containing
$\mathcal{C}$.
\end{proposition}

\begin{remark*}[Fully quantified form]
$\forall\mathcal{C}\subseteq\mathcal{P}(X)$:
(i) $\sigma(\mathcal{C})$ is a $\sigma$-algebra on $X$,
(ii) $\mathcal{C}\subseteq\sigma(\mathcal{C})$,
(iii) $\forall\mathcal{F}\ \sigma\text{-algebra on }X$ with
$\mathcal{C}\subseteq\mathcal{F}$:
$\sigma(\mathcal{C})\subseteq\mathcal{F}$.
\end{remark*}

\begin{remark*}[Flash]
\FlashQ{Outline the proof that $\sigma(\mathcal{C})$ exists.}
\FlashA{Non-emptiness: $\mathcal{P}(X)$ is a $\sigma$-algebra containing
  $\mathcal{C}$, so the intersecting family is non-empty.
  $\sigma(\mathcal{C})$ is a $\sigma$-algebra: by
  Proposition~\ref{prop:intersection-sigma-algebras}, arbitrary intersections
  of $\sigma$-algebras are $\sigma$-algebras.
  Minimality: if $\mathcal{F}$ contains $\mathcal{C}$ then $\mathcal{F}$ is in
  the intersecting family, so $\sigma(\mathcal{C})\subseteq\mathcal{F}$.}
\FlashHint{The proof is essentially one sentence once you have
  Proposition~\ref{prop:intersection-sigma-algebras}.}
\FlashImplies{prop:generated-sigma-idempotent, prop:generated-sigma-monotone}
\end{remark*}

\begin{proposition}[\texorpdfstring{%
  \hyperref[prf:generated-sigma-idempotent]{%
    Generated $\sigma$-Algebra is Idempotent}}{%
  Generated Sigma-Algebra is Idempotent}]
\label{prop:generated-sigma-idempotent}
For any $\mathcal{C} \subseteq \mathcal{P}(X)$:
$\sigma(\sigma(\mathcal{C})) = \sigma(\mathcal{C})$.
\end{proposition}

\begin{remark*}[Fully quantified form]
$\forall\mathcal{C}\subseteq\mathcal{P}(X):\
\sigma(\sigma(\mathcal{C})) = \sigma(\mathcal{C})$.
\end{remark*}

\begin{remark*}[Flash]
\FlashQ{Prove $\sigma(\sigma(\mathcal{C})) = \sigma(\mathcal{C})$.}
\FlashA{($\supseteq$): $\sigma(\mathcal{C})\subseteq\sigma(\sigma(\mathcal{C}))$
  since $\sigma(\mathcal{C})\subseteq\sigma(\mathcal{C})$ (trivially $\mathcal{C}$
  generates a set containing itself).
  ($\subseteq$): $\sigma(\mathcal{C})$ is itself a $\sigma$-algebra containing
  $\sigma(\mathcal{C})$, so by minimality
  $\sigma(\sigma(\mathcal{C}))\subseteq\sigma(\mathcal{C})$.}
\FlashHint{The key is that $\sigma(\mathcal{C})$ is already a $\sigma$-algebra,
  so generating from it again gives nothing new.}
\FlashImplies{prop:generated-sigma-monotone}
\end{remark*}

\begin{proposition}[\texorpdfstring{%
  \hyperref[prf:generated-sigma-monotone]{%
    Generated $\sigma$-Algebra is Monotone}}{%
  Generated Sigma-Algebra is Monotone}]
\label{prop:generated-sigma-monotone}
If $\mathcal{C}_1 \subseteq \mathcal{C}_2 \subseteq \mathcal{P}(X)$
then $\sigma(\mathcal{C}_1) \subseteq \sigma(\mathcal{C}_2)$.
\end{proposition}

\begin{remark*}[Fully quantified form]
$\forall\mathcal{C}_1,\mathcal{C}_2\subseteq\mathcal{P}(X)$:
$\mathcal{C}_1\subseteq\mathcal{C}_2 \Rightarrow
\sigma(\mathcal{C}_1)\subseteq\sigma(\mathcal{C}_2)$.
\end{remark*}

\begin{remark*}[Flash]
\FlashQ{Prove monotonicity of the generated $\sigma$-algebra.}
\FlashA{$\sigma(\mathcal{C}_2)$ is a $\sigma$-algebra containing
  $\mathcal{C}_2\supseteq\mathcal{C}_1$.
  So $\sigma(\mathcal{C}_2)$ is a $\sigma$-algebra containing $\mathcal{C}_1$.
  By minimality of $\sigma(\mathcal{C}_1)$:
  $\sigma(\mathcal{C}_1)\subseteq\sigma(\mathcal{C}_2)$.}
\FlashHint{One line once you invoke minimality correctly.}
\FlashImplies{def:generated-sigma-algebra}
\end{remark*}

\begin{proposition}[\texorpdfstring{%
  \hyperref[prf:countable-cocountable-sigma]{%
    Countable/Co-Countable Sets Form a $\sigma$-Algebra}}{%
  Countable Co-Countable Sigma-Algebra}]
\label{prop:countable-cocountable-sigma}
Let $X$ be any non-empty set.  The collection
\[
  \mathcal{F} := \{ A \subseteq X :
    A \text{ is countable or } A^c \text{ is countable} \}
\]
is a $\sigma$-algebra on $X$.
\end{proposition}

\begin{remark*}[Fully quantified form]
$\forall$ non-empty set $X$:
$\mathcal{F} = \{A\subseteq X : |A|\leq\aleph_0\text{ or }|A^c|\leq\aleph_0\}$
is a $\sigma$-algebra on $X$.
\end{remark*}

\begin{remark*}[Flash]
\FlashQ{Verify the three $\sigma$-algebra axioms for the countable/co-countable
  collection.}
\FlashA{($\sigma$1): $\emptyset$ is countable, so $\emptyset\in\mathcal{F}$.
  ($\sigma$2): if $A$ is countable then $A^c$ has countable complement $A$,
  so $A^c\in\mathcal{F}$; if $A^c$ is countable then $A\in\mathcal{F}$
  trivially.
  ($\sigma$3): if all $A_n$ are countable then $\bigcup A_n$ is countable
  (countable union of countable sets); if some $A_k^c$ is countable then
  $(\bigcup A_n)^c\subseteq A_k^c$ is countable.}
\FlashHint{For ($\sigma$3) split into two cases: all $A_n$ countable, or
  at least one $A_k$ co-countable.  The second case uses monotonicity of
  complement.}
\FlashImplies{prop:finite-cocountable-algebra-not-sigma}
\end{remark*}

\begin{proposition}[\texorpdfstring{%
  \hyperref[prf:finite-cocountable-algebra-not-sigma]{%
    Finite/Co-Finite Sets are an Algebra but not a $\sigma$-Algebra}}{%
  Finite Co-Finite Algebra Not Sigma-Algebra}]
\label{prop:finite-cocountable-algebra-not-sigma}
The collection
$\mathcal{A} := \{ A \subseteq \mathbb{N} :
  A \text{ is finite or } A^c \text{ is finite} \}$
is an algebra of sets on $\mathbb{N}$ but is not a $\sigma$-algebra
on $\mathbb{N}$.
\end{proposition}

\begin{remark*}[Fully quantified form]
$\mathcal{A} = \{A\subseteq\mathbb{N}: |A|<\infty\text{ or }|A^c|<\infty\}$
satisfies (A1)--(A3) but fails ($\sigma$3).
\end{remark*}

\begin{remark*}[Flash]
\FlashQ{Prove that the finite/co-finite algebra on $\mathbb{N}$ is not a
  $\sigma$-algebra.}
\FlashA{It is an algebra: $\emptyset$ finite; complement swaps finite and co-finite;
  union of two finite sets is finite; union of a finite and co-finite set is co-finite.
  It fails ($\sigma$3): for each $n$, $\{2n\}\in\mathcal{A}$ (singleton, finite).
  But $\bigcup_{n=0}^\infty\{2n\} = \{0,2,4,\ldots\}$ is infinite and its
  complement $\{1,3,5,\ldots\}$ is also infinite, so
  $\{0,2,4,\ldots\}\notin\mathcal{A}$.}
\FlashHint{The witness is any infinite co-infinite subset of $\mathbb{N}$
  expressed as a countable union of singletons.}
\FlashImplies{def:sigma-algebra, prop:hierarchy-inclusions}
\end{remark*}

\begin{proposition}[\texorpdfstring{%
  \hyperref[prf:sigma-from-singletons]{%
    Singletons Generate the Full Power Set on Countable Sets}}{%
  Singletons Generate Power Set on Countable Sets}]
\label{prop:sigma-from-singletons}
If $X$ is a countable non-empty set, then
$\sigma\bigl(\{\{x\} : x \in X\}\bigr) = \mathcal{P}(X)$.
\end{proposition}

\begin{remark*}[Fully quantified form]
$\forall$ countable non-empty set $X$:
$\sigma(\{\{x\}:x\in X\}) = \mathcal{P}(X)$.
\end{remark*}

\begin{remark*}[Flash]
\FlashQ{Why does the $\sigma$-algebra generated by singletons equal
  $\mathcal{P}(X)$ when $X$ is countable?}
\FlashA{($\subseteq$): $\sigma(\{\{x\}:x\in X\})\subseteq\mathcal{P}(X)$
  trivially.
  ($\supseteq$): let $A\subseteq X$; since $X$ is countable, $A$ is countable,
  so $A = \bigcup_{x\in A}\{x\}$ is a countable union of singletons, each in
  $\sigma(\{\{x\}:x\in X\})$, hence $A\in\sigma(\{\{x\}:x\in X\})$.}
\FlashHint{Countability of $A$ is essential --- on uncountable sets singletons
  generate only the countable/co-countable $\sigma$-algebra.}
\FlashImplies{prop:countable-cocountable-sigma}
\end{remark*}

\begin{proposition}[\texorpdfstring{%
  \hyperref[prf:sigma-closed-symmdiff]{%
    $\sigma$-Algebra Closed Under Countable Symmetric Difference}}{%
  Sigma-Algebra Closed Under Countable Symmetric Difference}]
\label{prop:sigma-closed-symmdiff}
Let $\mathcal{F}$ be a $\sigma$-algebra on $X$ and let
$(A_n)_{n\in\mathbb{N}} \subseteq \mathcal{F}$.  Define
$A_n \mathbin{\triangle} B_n$ for any $(B_n)\subseteq\mathcal{F}$.
Then $\mathcal{F}$ is closed under pairwise symmetric difference:
$A, B \in \mathcal{F} \Rightarrow A \mathbin{\triangle} B \in \mathcal{F}$.
\end{proposition}

\begin{remark*}[Fully quantified form]
$\forall\mathcal{F}\ \sigma\text{-algebra on }X$,
$\forall A,B\in\mathcal{F}:\
A\triangle B\in\mathcal{F}$.
\end{remark*}

\begin{remark*}[Flash]
\FlashQ{Why is every $\sigma$-algebra closed under symmetric difference?}
\FlashA{$A\triangle B = (A\setminus B)\cup(B\setminus A)
  = (A\cap B^c)\cup(B\cap A^c)$.
  Each of $B^c, A^c\in\mathcal{F}$ by complement closure;
  each intersection $\in\mathcal{F}$ by finite intersection closure
  (which follows from complement and union);
  the union $\in\mathcal{F}$ by union closure.}
\FlashHint{Express $\triangle$ entirely in terms of $\cup$, $\cap$, $(\cdot)^c$
  and verify each step uses a $\sigma$-algebra axiom or a derived closure property.}
\FlashImplies{thm:powerset-boolean-ring}
\end{remark*}
